\SourceLeaf{386}%
die Messungen selbst erlauben, sodass man aus jenen (den geodätischen Längen und Breiten) die gemessenen Winkel rückwärts wenigstens mit derselben Genauigkeit berechnen kann, mit der sie gemessen wurden, so reichen die abgekürzten Formeln bei Weitem nicht aus. Bei sehr großen Dreiecken wäre es dann sogar wünschenswert, auch noch die Glieder vierter Ordnung berücksichtigen zu können. Bei einer anderen Form der Rechnung lässt sich das mit sehr flexiblen Methoden erreichen. Hier ist jedoch nicht der Ort, das weiter auszuführen, und wir begnügen uns damit, auf dieses Bedürfnis der höheren Geodäsie hinzuweisen. Das Werk selbst bietet verschiedene Fälle, in denen die großen Vorteile einer solchen Behandlung spürbar werden. So wurden z.\,B. die bei den Azimutbestimmungen an der Turiner Sternwarte vorgenommenen Einschneidungen der Dreieckspunkte Masse, Monte Soglio und Rocca Melone überhaupt nicht genutzt. Unter der genannten Voraussetzung hätten sie ohne Mühe eine sehr wertvolle Kontrolle und eine größere Genauigkeit ermöglicht.

